
\chapter{DevOps}
\label{ch:DevOps} % Label for method chapter
\section{Semantic Versioning e Changelog e secrets}
\subsection{Semantic Versioning}
Per ogni release è stato adottato il semanting Versioning. È stato utilizzato il formato X.Y.Z, partendo da 1.0.0, ogni release incrementa il numero di versione in base a:
\begin{itemize}
    \item \textbf{Major (X):} per cambiamenti incompatibili con le versioni precedenti
    \item \textbf{Minor (Y):} per aggiunta di nuove funzionalità compatibili con le versioni precedenti
    \item \textbf{Patch (Z):} per correzione di bug compatibili con le versioni precedenti
\end{itemize}
Per ogni release viene creato anche un tag con la corrispondente versione.
\subsection{Changelog e secrets}
Ogni volta che is effettua una push sul main, viene generato un changelog. Il changelog è un file che contiene tutte le modifiche effettuate.
Nel changelog vengono riportati ogni commit che indicano le modifiche apportate, i contributori e il codice sorgente.
Se presenti, vengono inoltre riportate le sezioni: 
\begin{itemize}
    \item \textbf{feat: Feature} per le nuove funzionalità
    \item \textbf{fix: Bug Fixes} per le correzioni di bug
    \item \textbf{docs: Documentation} per le modifiche alla documentazione
    \item \textbf{refactor: Refactor} per le modifiche al codice che non aggiungono nuove funzionalità
    \item \textbf{perf: Performance Improvements} per le modifiche che migliorano le prestazioni
    \item \textbf{CI: Continuos Integration} per le modifiche alla CI
\end{itemize}
Il changelog viene autogenerato tramite l'action 'BobAnkh/auto-generate-changelog@v1.2.5'.
Inoltre per evitare un accesso non autorizzato a dati sensibili, come l'access token del changelog, sono stati inseriti in secrets, variabili d'ambiente di Github. 
\section{Licensing e issue template}
In ogni nostra repo è presente una licenza MIT. È stata scelta questa licenza perché è estremamente permissiva: 
il codice sorgente di Maraffa Online è open source e pertanto gli utenti possono utilizzarlo, modificarlo e distribuirlo.
Inoltre la licenza MIT è compatibile anche con altre licenze, anche con quelle restrittive come GPL.
È incoraggiata la collaborazione anche grazie alla presenza di un template con il quale un utente può consigliare nuove feature,
 segnalare un bug o suggerire un'implementazione alternativa.
\section{Git section; non so come chiamarla}
\subsection{Strategy for version control}
**parla della git policy**
\subsection{Git workflow}
**main, develop, feature branch...***
\subsection{Commit verification}
È stato utilizzato il meccanismo di commit verification per assicurare:
\begin{itemize}
    \item \textbf{Integrità:} controllo sull'alterazione dei file presenti all'interno del commit.  Utilizzando una funzione hash, ci si assicura che non siano stati alterati in non modo autorizzato.
    \item \textbf{Autenticità:} controllo sull'autore del commit. Utilizzando una firma digitale, ci si assicura che il commit sia stato effettuato dall'autore autorizzato.
\end{itemize}
I controlli vengono effettuati in seguito a un merge. \\
Quando un commit è verificato viene prodotto un tag verde con la scritta "verified" e un tag rosso con la scritta "unverified" in caso contrario:
\begin{figure}[h!]
    \centering 
    \includegraphics[scale=0.7]{report/img/verified_commit.png}
    \caption{Verified commit}
    \label{verified_commit}
    \end{figure}
\subsection{Conventional commits}
In ogni repo è stato adottato un sistema standard per scrivere i commit:  i conventional commit. In questo modo i commit risultano più chiari e facilmente leggibili. 
Riportiamo di seguito la nomenclatura usata:
\begin{itemize}
    \item \textbf{fix:} per i commit che risolvono un bug
    \item \textbf{feat:} per i commit che aggiungono una nuova feature
    \item \textbf{refactor:} per i commit che migliorano il codice senza aggiungere nuove funzionalità
    \item \textbf{docs:} per i commit che riguardano la documentazione
    \item \textbf{style:} per i commit che riguardano la formattazione del codice
    \item \textbf{test:} per i commit che riguardano i test
    \item \textbf{ci:} per i commit che riguardano la Continuous Integration
\end{itemize}
Inoltre sono evidenziati i breaking changes per le modifiche non più compatibili con le versioni precedenti.
La stessa nomenclatura viene utilizzata anche nel changelog.

\section{Build automation}
\section{Continuous Integration}
Abbiamo deciso di implementare ogni servizio in una repository diversa. La repository fornita è un contenitore di tutte le singole repository dei servizi.
Per ogni repository è stata implementata una CI specifica. \\

\textbf{UserManagementMaraffa e BusinessLogic:} Essendo le due repository estremamente simili e avendo entrambe un ambiente in node.js, è stato implementato un workflow pressochè analogo. Questo workflow reagisce all'evento push sul main branch, e pull request sui branch main e develop 
e sono stati implementati 3 job: build, test e deploy.
Dopo una checkout iniziale (fetch-depth è 0 in quanto il changelog necessita di tutta la history dei commit), viene creata una cache per permettere ai job di comunicare tra loro. Successivamente alla build il sistema viene testato (in entrambi i casi viene usata l'action di yarn) e infine viene effettuato il deploy del servizio containerizzato se l'evento è una pull request con la quale viene svolto il merge e quando l'action è closed (non sta progredendo e quindi l'esecuzione non è in corso). Nel deploy viene effettuato il chekout, si configura il nome del progetto, si configura la versione della build. Se il commit viene effettuato sul main si configura il nome del tag e la versione a "LATEST", nel caso non fosse nel main viene configurata la versione a "develop". Infine si effettua il login al Github container registry, la build e la push dell'immagine sul registro. L'immagine viene configurata in modo parametrico con delle variabili d'ambiente e il nome viene trasformato in minuscolo. 

\textbf{MiddlewareMaraffa:} //TODO
push su main, pull request su main e develop
build e test se pullrequest.merge = true 

Parte finale 
\textbf{ChatServer:} Il servizio è sviluppato in un ambiente Java, pertanto dopo aver effettuato il checkout nel job build and test,  si effettua il set-up di Java (distribuzione: temurin, versione 20). Viene svolto il set-up di Gradle, in cui viene usata un'action per la build e creata una cache. Inoltre viene aggiunto un commento con un riassunto dei risultati dei job  alla pull request, il cui valore è impostato a "on-failure". Si noti che il suo valore di default è "never"; se fosse "always", si potrebbero verificare scenari indesiderati. Il riassunto dei risultati dei job risulta particolarmente utile: un utente che vuole esaminare la pull request riesce a comprendere lo stato dei job (se sono stati eseguiti correttamente oppure no) in modo leggibile e intuitivo. Successivamente si effettua il deploy del container di rabbitMQ se l'evento è una pull request oppure se è una push sul master o se la stringa con il riferimento a github contiene "tags/v". Infine se tutti il job da cui dipende (build-and-test) non è fallito e non è stato interrotto. Con l'istruzione run si verifica che non ci siano stati dei fallimenti nel job deploy and test.

\textbf{FrontendMaraffa:} //TODO
Sarebbe figo spiegare come i nostri workflow hanno ridotto i conflitti di integrazione, il rilevamento precoce di bug e i problemi di compatibilità.

\section{Automated quality assurance}
**code quality, testing, reporting**
\section{Automated delivery}
**automatizzare l'esecuzione iniziale del sistema con build.gradle che specifica come
il file deve compilare, gestire dip e configurare altri apetti del build process**
\section{Deployment}
è opzionale
\section{Containerization}
** docker image for each microservice, docker file
docker compose grouping containers and defining dependencies**

% We mentioned in Chapter~\ref{ch:into} %[example backward reference to a chapter or section.]
% that a project report's structure could follow a particular paradigm. Hence, the organization of a report (effectively the Table of Content of a report) can vary depending on the type of project you are doing. Check which of the given examples suit your project. Alternatively, follow your supervisor's advice.

% \section{Examples of the sections of a methodology chapter}
% A general report structure is summarised (suggested) in Table~\ref{tab:gen_template}. Table~\ref{tab:gen_template} describes that, in general, a typical report structure has three main parts: (1) front matter, (2) main text, and (3) end matter. %[\textbf{also notice that the preceding sentence is an example of a numbered list in a text body}]. 
% The structure of the front matter and end matter will remain the same for all the undergraduate final year project report. However, the main text varies as per the project's needs.
% \begin{table}[h!]
%     \centering
%     \caption{Undergraduate report template structure}
%     \label{tab:gen_template}
%     \begin{tabular}{llll}     
%         \toprule
%         \multirow{7}{3cm}{Frontmatter} 
%         & & Title Page & \\                  
%         & & Abstract &    \\          
%         & & Acknowledgements & \\                            
%         & & Table of Contents &    \\                                
%         & & List of Figures   &    \\                        
%         & & List of Tables    &    \\                
%         & & List of Abbreviations  &    \\                     
%         & &   &    \\                        
%         \multirow{7}{3cm}{Main text}
%         & Chapter 1 & Introduction   &    \\                         
%         & Chapter 2 & Literature Review   &    \\
%         & Chapter 3 & Methodology   &    \\
%         & Chapter 4 & Results    &    \\
%         & Chapter 5 & Discussion and Analysis  &    \\
%         & Chapter 6 & Conclusions and Future Work  &    \\        
%         & Chapter 7 & Refection  &    \\          
%         & &   &    \\                       
%         \multirow{2}{3cm}{End matter}
%         & & References  &    \\   
%         & & Appendices (Optional)  &    \\ 
%         & & Index (Optional)  &    \\ 
%         \bottomrule
%     \end{tabular}
% \end{table}

% \subsection{Example of a software/Web development main text structure}
% \label{subsec:se_chpters}
% Notice that the ``methodology'' Chapter of Software/Web development in Table~\ref{tab:soft_eng_temp} takes a standard software engineering paradigm (approach). Alternatively, these suggested sections can be the chapters of their own. Also, notice that ``Chapter 5'' in Table~\ref{tab:soft_eng_temp} is ``Testing and Validation'' which is different from the general report template mentioned in Table~\ref{tab:gen_template}. Check with your supervisor if in doubt.
% \begin{table}[h!]
%     \centering
%     \caption{Example of a software engineering-type report structure}
%     \label{tab:soft_eng_temp}
%     \begin{tabular}{lll}     
%         \toprule                   
%         Chapter 1 & Introduction   &    \\        
%         Chapter 2 & Literature Review  &    \\                   
%         Chapter 3 & Methodology   &    \\
%         &               & Requirements specifications   \\
%         &               & Analysis   \\
%         &               & Design   \\
%         &               & Implementations   \\
%         Chapter 4 & Testing and Validation  &    \\
%         Chapter 5 & Results and Discussion      &    \\
%         Chapter 6 & Conclusions and Future Work  &    \\        
%         Chapter 7 & Reflection  &    \\                          
%         \bottomrule
%     \end{tabular}
% \end{table}

% \subsection{Example of an algorithm analysis main text structure}
% Some project might involve the implementation of a state-of-the-art algorithm and its performance analysis and comparison with other algorithms. In that case, the suggestion in Table~\ref{tab:algo_temp} may suit you the best. 
% \begin{table}[h!]
%     \centering
%     \caption{Example of an algorithm analysis type report structure}
%     \label{tab:algo_temp}
%     \begin{tabular}{lll}     
%         \toprule                   
%         Chapter 1 & Introduction  &    \\        
%         Chapter 2 & Literature Review  &    \\                
%         Chapter 3 & Methodology   &    \\
%         &               & Algorithms descriptions  \\
%         &               & Implementations   \\
%         &               & Experiments design   \\
%         Chapter 4 & Results       &  \\
%         Chapter 5 & Discussion and Analysis  &    \\
%         Chapter 6 & Conclusion and Future Work  &    \\        
%         Chapter 7 & Reflection  &    \\          
%         \bottomrule
%     \end{tabular}
% \end{table}

% \subsection{Example of an application type main text structure}
% If you are applying some algorithms/tools/technologies on some problems/datasets/etc., you may use the methodology section prescribed in Table~\ref{tab:app_temp}.  
% \begin{table}[h!]
%     \centering
%     \caption{Example of an application type report structure}
%     \label{tab:app_temp}
%     \begin{tabular}{lll}     
%         \toprule                   
%         Chapter 1 & Introduction  &    \\        
%         Chapter 2 & Literature Review  &    \\                
%         Chapter 3 & Methodology   &    \\
%         &               & Problems (tasks) descriptions  \\
%         &               & Algorithms/tools/technologies/etc. descriptions  \\        
%         &               & Implementations   \\
%         &               & Experiments design and setup   \\
%         Chapter 4 & Results       &  \\
%         Chapter 5 & Discussion and Analysis  &    \\
%         Chapter 6 & Conclusion and Future Work  &    \\        
%         Chapter 7 & Reflection  &    \\          
%         \bottomrule
%     \end{tabular}
% \end{table}

% \subsection{Example of a science lab-type main text structure}
% If you are doing a science lab experiment type of project, you may use the  methodology section suggested in Table~\ref{tab:lab_temp}. In this kind of project, you may refer to the ``Methodology'' section as ``Materials and Methods.''
% \begin{table}[h!]
%     \centering
%     \caption{Example of a science lab experiment-type report structure}
%     \label{tab:lab_temp}
%     \begin{tabular}{lll}     
%         \toprule                   
%         Chapter 1 & Introduction  &    \\        
%         Chapter 2 & Literature Review  &    \\                
%         Chapter 3 & Materials and Methods   &    \\
%         &               & Problems (tasks) description  \\
%         &               & Materials \\        
%         &               & Procedures  \\                
%         &               & Implementations   \\
%         &               & Experiment set-up   \\
%         Chapter 4 & Results       &  \\
%         Chapter 5 & Discussion and Analysis  &    \\
%         Chapter 6 & Conclusion and Future Work  &    \\        
%         Chapter 7 & Reflection  &    \\          
%         \bottomrule
%     \end{tabular}
% \end{table}

% \section{Example of an Equation in \LaTeX}
% Eq.~\ref{eq:eq_example} [note that this is an example of an equation's in-text citation] is an example of an equation in \LaTeX. In Eq.~\eqref{eq:eq_example}, $ s $ is the mean of elements $ x_i \in \mathbf{x} $: 

% \begin{equation}
% \label{eq:eq_example} % label used to refer the eq in text
% s = \frac{1}{N} \sum_{i = 1}^{N} x_i. 
% \end{equation}

% Have you noticed that all the variables of the equation are defined using the \textbf{in-text} maths command \$.\$, and Eq.~\eqref{eq:eq_example} is treated as a part of the sentence with proper punctuation? Always treat an equation or expression as a part of the sentence. 

% \section{Example of a Figure in \LaTeX}
% Figure~\ref{fig:chart_a} is an example of a figure in \LaTeX. For more details, check the link:

% \href{https://en.wikibooks.org/wiki/LaTeX/Floats,_Figures_and_Captions}{wikibooks.org/wiki/LaTeX/Floats,\_Figures\_and\_Captions}.

% \noindent
% Keep your artwork (graphics, figures, illustrations) clean and readable. At least 300dpi is a good resolution of a PNG format artwork. However, an SVG format artwork saved as a PDF will produce the best quality graphics. There are numerous tools out there that can produce vector graphics and let you save that as an SVG file and/or as a PDF file. One example of such a tool is the ``Flow algorithm software''. Here is the link for that: \href{http://www.flowgorithm.org/download/}{flowgorithm.org}.
% \begin{figure}[ht]
%     \centering
%     % \includegraphics[scale=0.3]{figures/chart.pdf}
%     \caption{Example figure in \LaTeX.}
%     \label{fig:chart_a}
% \end{figure}

% \clearpage %  use command \clearpage when you want section or text to appear in the next page.

% \section{Example of an algorithm in \LaTeX}
% Algorithm~\ref{algo:algo_example} is a good example of an algorithm in \LaTeX.  
% \begin{algorithm}
%     \caption{Example caption: sum of all even numbers}
%     \label{algo:algo_example}
%     \begin{algorithmic}[1]
%         \Require{$ \mathbf{x}  = x_1, x_2, \ldots, x_N$}
%         \Ensure{$EvenSum$ (Sum of even numbers in $ \mathbf{x} $)}
%         \Statex
%         \Function{EvenSummation}{$\mathbf{x}$}
%         \State {$EvenSum$ $\gets$ {$0$}}
%         \State {$N$ $\gets$ {$length(\mathbf{x})$}}
%         \For{$i \gets 1$ to $N$}                    
%         \If{$ x_i\mod 2 == 0$}  \Comment check if a number is even?
%         \State {$EvenSum$ $\gets$ {$EvenSum + x_i$}}
%         \EndIf
%         \EndFor
%         \State \Return {$EvenSum$}
%         \EndFunction
%     \end{algorithmic}
% \end{algorithm}
 
% \section{Example of code snippet  in \LaTeX}

% Code Listing~\ref{list:python_code_ex} is a good example of including a code snippet in a report. While using code snippets, take care of the following:
% \begin{itemize}
%     \item do not paste your entire code (implementation) or everything you have coded. Add code snippets only. 
%     \item The algorithm shown in Algorithm~\ref{algo:algo_example} is usually preferred over code snippets in a technical/scientific report. 
%     \item Make sure the entire code snippet or algorithm stays on a single page and does not overflow to another page(s).  
% \end{itemize}

% Here are three examples of code snippets for three different languages (Python, Java, and CPP) illustrated in Listings~\ref{list:python_code_ex}, \ref{list:java_code_ex}, and \ref{list:cpp_code_ex} respectively.  

% \begin{lstlisting}[language=Python, caption={Code snippet in \LaTeX ~and  this is a Python code example}, label=list:python_code_ex]
% import numpy as np

% x  = [0, 1, 2, 3, 4, 5] # assign values to an array
% evenSum = evenSummation(x) # call a function

% def evenSummation(x):
%     evenSum = 0
%     n = len(x)
%     for i in range(n):
%         if np.mod(x[i],2) == 0: # check if a number is even?
%             evenSum = evenSum + x[i]
%     return evenSum
% \end{lstlisting}

% Here we used  the ``\textbackslash clearpage'' command and forced-out the second listing example onto the next page. 
% \clearpage  %
% \begin{lstlisting}[language=Java, caption={Code snippet in \LaTeX ~and  this is a Java code example}, label=list:java_code_ex]
% public class EvenSum{ 
%     public static int evenSummation(int[] x){
%         int evenSum = 0;
%         int n = x.length;
%         for(int i = 0; i < n; i++){
%             if(x[i]%2 == 0){ // check if a number is even?
%                 evenSum = evenSum + x[i];
%             }
%         }
%         return evenSum;     
%     }
%     public static void main(String[] args){ 
%         int[] x  = {0, 1, 2, 3, 4, 5}; // assign values to an array
%         int evenSum = evenSummation(x);
%         System.out.println(evenSum);
%     } 
% } 
% \end{lstlisting}


% \begin{lstlisting}[language=C, caption={Code snippet in \LaTeX ~and  this is a C/C++ code example}, label=list:cpp_code_ex]
% int evenSummation(int x[]){
%     int evenSum = 0;
%     int n = sizeof(x);
%     for(int i = 0; i < n; i++){
%         if(x[i]%2 == 0){ // check if a number is even?
%             evenSum = evenSum + x[i];
%     	}
%     }
%     return evenSum;     
% }

% int main(){
%     int x[]  = {0, 1, 2, 3, 4, 5}; // assign values to an array
%     int evenSum = evenSummation(x);
%     cout<<evenSum;
%     return 0;
% }
% \end{lstlisting}



% \section{Example of in-text citation style}
% \subsection{Example of the equations and illustrations placement and reference in the text}
% Make sure whenever you refer to the equations, tables, figures, algorithms,  and listings for the first time, they also appear (placed) somewhere on the same page or in the following page(s). Always make sure to refer to the equations, tables and figures used in the report. Do not leave them without an \textbf{in-text citation}. You can refer to equations, tables and figures more them once.

% \subsection{Example of the equations and illustrations style}
% Write \textbf{Eq.} with an uppercase ``Eq`` for an equation before using an equation number with (\textbackslash eqref\{.\}). Use ``Table'' to refer to a table, ``Figure'' to refer to a figure, ``Algorithm'' to refer to an algorithm and ``Listing'' to refer to listings (code snippets). Note that, we do not use the articles ``a,'' ``an,'' and ``the'' before the words Eq., Figure, Table, and Listing, but you may use an article for referring the words figure, table, etc. in general.

% For example, the sentence ``A report structure is shown in \textbf{the} Table~\ref{tab:gen_template}'' should be written as ``A report structure is shown \textbf{in} Table~\ref{tab:gen_template}.'' 
 

% \section{Summary}
% Write a summary of this chapter.

% ~\\[5em]
% \noindent
% {\huge\textbf{Note:}} In the case of \textbf{software engineering} project a Chapter ``\textbf{Testing and Validation}'' should precede the ``Results'' chapter. See Section~\ref{subsec:se_chpters} for report organization of such project. 

